\section{Protected integration and orientation attack repair}
\label{sec:integration-orientation-attack-repair}

\subsection{Claim repaired}

The repaired claim is a triad.

\begin{enumerate}
\item The reduced OP/DT/PT theorem is scalar.  Its integration map is the
Behrend-weighted reduced quotient Euler characteristic on
\(\mathrm{Hilb}^n(X,(\beta,d))/E\) and \(P_n(X,(\beta,d))/E\), and the
proved branch is
\[
  Z^X_{\mathrm{DT}}(T,Q,P)=Z^X_{\mathrm{OP}}(P,Q,T)
  =-4096\,\Delta_5(Z)^{-2}.
\]
\item A protected Hall or critical CoHA integration map is only a
criterion at this stage.  It requires oriented critical extension
correspondences, vanishing-cycle coefficients, HN-completed sector
quotients, wall-crossing compatibility, primitive extraction, and a
lifted Weyl action on the oriented completed theory.
\item The orientation character is an open geometric problem.  One must
construct the monodromy character
\[
  \epsilon_o:W^{(2)}(\Lambda^{2,1}_{II})\to\{\pm1\}
\]
from the square root of the Hall determinant line and prove
\[
  \epsilon_o|_{W^{(2)}(\Lambda^{2,1}_{II})}
  =
  \nu_{\Delta_5}|_{W^{(2)}(\Lambda^{2,1}_{II})}.
\]
This is not produced by normalizing \(\Delta_5\).
\end{enumerate}

\subsection{Constants separated}

The manuscript now separates the constants and signs.

\[
  [q^{1/2}r^{1/2}s^{1/2}]\Delta_5=64
\]
is the theta-normalized leading coefficient, equivalently the coefficient
attached to the Weyl monomial in the determinant.

\[
  D_5=64^{-1}\Delta_5,\qquad
  \chi_{10}^{\mathrm{OP}}=D_5^2=4096^{-1}\Delta_5^2
\]
is the OP monic normalization.  Thus \(4096=64^2\).

\[
  Z^X_{\mathrm{OP}}=-(\chi_{10}^{\mathrm{OP}})^{-1}
\]
is the OP scalar sign convention.  None of \(64\), \(4096\), or the OP
minus sign is the Hall orientation character.

\subsection{Files changed}

\begin{enumerate}
\item \(\texttt{main.tex}\): rewrote the owned
\emph{The protected-index dictionary} and \emph{The scalar square}
sections into the theorem/criterion/open-problem triad; updated only the
orientation item in Definition~\ref{def:compact-igusa-realization-datum}
to state that \(\epsilon_o\) is geometric monodromy, not scalar
normalization.
\item \(\texttt{appendices/boundary\_compatibility\_conditions.tex}\):
updated the orientation-character row and paragraph to say explicitly
that the comparison character must come from Hall determinant-line
monodromy.
\item \(\texttt{agent\_material/13\_integration\_orientation\_attack\_repair.tex}\):
this report.
\end{enumerate}

\subsection{Checks run}

\begin{enumerate}
\item Read the required files:
\(\texttt{CLAUDE.md}\), \(\texttt{AGENTS.md}\),
\(\texttt{\string~/ecosystem/INVARIANTS.md}\), \(\texttt{main.tex}\),
\(\texttt{agent\_material/12\_protected\_integration\_orientation\_attack\_heal.tex}\),
and \(\texttt{agent\_material/03\_scalar\_square\_status.tex}\).
\item Checked bibliography keys for
\(\texttt{OB}\), \(\texttt{OPand}\),
\(\texttt{OberdieckReducedPT}\), \(\texttt{BRYAN}\),
\(\texttt{BEHREND}\), \(\texttt{DT}\),
\(\texttt{KSCoHA}\), and \(\texttt{DavisonMeinhardtCoHA}\).
\item Ran targeted \(\texttt{rg}\) scans for
\(\epsilon_o\), \(\nu_{\Delta_5}\), \(64\), \(4096\), protected
integration, Hall orientation, HN completion, vanishing cycles, and wall
crossing.
\item Ran \(\texttt{make}\).  It completed successfully and rebuilt
\(\texttt{out/main.pdf}\).
\item Scanned \(\texttt{out/main.log}\) for undefined references and
undefined citations; no such hits were reported.
\item Ran the invariant-language scan for attribution markers and forbidden
standalone-document phrases on the edited TeX surfaces; no hits.
\end{enumerate}

\subsection{Remaining obligations}

\begin{enumerate}
\item Construct the compact reduced \(K3\times E\) Hall or critical CoHA
sector with oriented extension correspondences.
\item Prove \(E\)-quotient descent for vanishing-cycle coefficients and
for the strong orientation datum.
\item Prove HN-finite completion and wall-crossing compatibility for
protected integration.
\item Construct primitive extraction and prove the Borcherds--Kac bracket
relations in the protected Hall/factorization bracket.
\item Compute the orientation-line monodromy for the simple diagonal
reflections \(s_{\delta_i}\), then extend it to
\[
  \epsilon_o|_{W^{(2)}(\Lambda^{2,1}_{II})}
  =
  \nu_{\Delta_5}|_{W^{(2)}(\Lambda^{2,1}_{II})}.
\]
\end{enumerate}
